\chapter{Implementation and Module-Wise Coding}
\label{ch:impl}

\section{Development Environment}
Python 3.12 (virtual environment), VS Code, Git and GitHub, Windows 11; Flask 3.1 with Jinja2, vanilla
JavaScript and Chart.js 4.4.1; scikit-learn 1.9.1, pandas 3.0, NumPy 2.5; SQLite 3 (WAL); DuckDB 1.5 for
reading the MIMIC-IV CSV files; pytest 9.1. AI coding assistants were used during development: 13 of the
25 non-merge commits carry an AI co-author trailer, and all such changes were reviewed and tested before
being committed.

\section{Overall Implementation Strategy}
The repository history shows the order of work: foundation and configuration, database schema and seeding,
feature pipeline, the four-model ensemble, the vault and QR engine, the REST API, the user interface and a
benchmark runner (6--7 October 2026); RBAC and the Gemini narrative (8 October); a security and RBAC fix
(14 October); then the research track---test isolation and data modes, the MIMIC-IV adapter, the Task~1
dataset, leakage-aware baselines, verified research serving and security fixes (25 October)---followed by
the research dashboard, model set v2, deployment hardening and a vitals-chart fix (26 October).

\section{Module-Wise Implementation}
\begin{small}
\setlength{\tabcolsep}{4pt}
\begin{longtable}{@{}p{24mm}p{66mm}p{40mm}p{28mm}@{}}
\caption{Module-wise implementation summary}\label{tab:impl}\\
\toprule
Module & Processing logic and key details & Code references & Test evidence\\
\midrule
\endfirsthead
\toprule Module & Processing logic and key details & Code references & Test evidence\\ \midrule
\endhead
Data pipeline & clean, validate, engineer 30 features, stratified 80/20 split, fit and save scaler & \code{src/data/*} & \code{test_data_pipeline.py}\\
Ensemble & train four models, fuse soft votes, rule path, KNN precedents, per-model scores & \code{src/models/*} & \code{test_models.py}, \code{test_benchmark.py}\\
API and RBAC & app factory, blueprints, session roles, error envelopes, rate limits, headers & \code{src/api/*} & \code{test_api.py}, \code{test_auth.py}\\
Vault and QR & HS256 JWT with DB record, scope filter, atomic quota, audit log & \code{src/vault/*} & \code{test_vault.py}\\
Narratives & de-identified prompt, TLS-verified call, local fallback & \code{src/nlp/summarizer.py} & \code{test_nlp.py}, \code{test_security_hardening.py}\\
MIMIC adapter and Task~1 & DuckDB reads, item-ID mapping, unit checks, cohort, labels, anomalies, hashed dataset & \code{src/datasources/mimic/*}, \code{src/tasks/*} & \code{test_mimic_t1_dataset.py}, \code{test_mimic_layout_mappings.py}\\
Research training & grouped CV, metrics, bootstrap, sensitivity analyses, staged atomic output & \code{src/training/*} & \code{test_training_t1.py}\\
Research serving & hash-verified loading, what-if contract, dashboard, out-of-fold stay view & \code{training/serving.py}, \texttt{stays.py}, \texttt{whatif.py} & \code{test_research_mode_app.py}, \code{test_whatif_contract.py}\\
Deployment & Waitress entry, first-start bootstrap, Docker image & \texttt{wsgi.py}, \code{src/bootstrap.py}, \texttt{Dockerfile} & \code{test_deployment_hardening.py}\\
\bottomrule
\end{longtable}
\end{small}

\section{Front-End Implementation}
Seven server-rendered portals share \texttt{base.html} (role-aware navigation, responsive header):
landing, login, physician triage (KPI cards, alert feed, filterable roster, chart drawer with Chart.js
vitals and model votes), patient vault (records, sharing dialog, QR image, access log), scanner (camera,
upload or pasted token), administrator overview, and in research mode the research triage and model
analytics pages. All scripts are static files: a shared \texttt{common.js} redirects to the login page on
an expired session, research scripts build the DOM with \texttt{textContent} only, and no data are kept in
browser storage. Screenshots appear in Chapter~\ref{ch:results}.

\section{Back-End Implementation}
The 44 routes (Appendix~\ref{app:repro}) follow one pattern: authenticate and authorise, validate the
JSON body (objects only, bounded values), call a service, and return a \texttt{\{success, data, error,
meta\}} envelope. Models are loaded once into application memory at start-up; research mode replaces them
with a \texttt{ResearchModelService} that verifies every artifact's SHA-256 hash before deserialisation
and blocks the synthetic clinical APIs with HTTP 503.

\section{Database Implementation}
\code{src/database/db.py} opens one connection per operation with \texttt{PRAGMA foreign\_keys=ON},
\texttt{journal\_mode=WAL} and \texttt{synchronous=NORMAL}, and returns dictionary-like rows; all SQL is
parameterised. Research mode opens the demo database read-only (\texttt{mode=ro}) and never creates it.

\section{Communication and Concurrency Implementation}
Clients communicate over HTTP with JSON; Waitress serves them with 8 worker threads. Shared state is
limited and protected: rate-limit counters and the research stay view use locks; alert acknowledgements are
kept per process, which is why a single instance is required. The redemption race (several concurrent
scans of a single-use pass all succeeding) was removed with an atomic conditional update and is covered by
a concurrency test.

\section{Git Repository Setup}
\label{sec:git}
\textbf{Repository:} \url{https://github.com/param7862989-ops/MediH} (30 commits; 144 tracked files;
about 24{,}400 lines of Python, JavaScript, HTML, CSS and SQL including tests).
\textbf{Branching:} \texttt{main} is the integration branch; work is done on personal feature branches
(\texttt{nirupam}, \texttt{param}) and merged through pull requests (five merged, \#1--\#5).
\textbf{Commit convention:} Conventional Commits (\texttt{feat(scope):}, \texttt{fix(scope):},
\texttt{docs:}). \textbf{History by Git identity:} \texttt{paramthakkar24-sudo} 13 commits,
\texttt{Nirupam} 10, \texttt{param7862989-ops} 7 (including the five merge commits). Git identities are not
verified mappings to students; individual contributions are to be confirmed by the group
(\textit{see MISSING\_INFORMATION.md}). At the time of writing, the research triage page (stay-level
out-of-fold view) is implemented and tested in the working tree of branch \texttt{param} but not yet
committed.

\section{Integration of Modules}
\texttt{create\_app()} integrates configuration, security checks, middleware, model loading and the
blueprints; \code{src/bootstrap.py} chains seeding $\rightarrow$ feature pipeline $\rightarrow$ ensemble
training when assets are missing; \code{run_demo.py} (local) and \texttt{wsgi.py} (production) are the two
entry points. The research pipeline is integrated only through its verified on-disk run.

\section{Implementation Challenges and Resolutions}
\begin{itemize}
\item \textbf{Look-ahead leakage:} a value can be charted inside the 24-hour window but stored in the
      record only later; the reader therefore requires \texttt{storetime} $\le$ index time as well
      (568 such observations excluded on the demo).
\item \textbf{Outcome timestamp anomalies:} four deaths are recorded after discharge; they are labelled
      positive, flagged and analysed separately instead of being silently dropped.
\item \textbf{Misleading what-if inputs:} measurement counts reflected monitoring intensity and could not
      be entered meaningfully; model set v2 excluded them (74 inputs).
\item \textbf{Concurrency defect:} single-use QR passes could be redeemed more than once; fixed with an
      atomic update.
\item \textbf{Stale demo data:} the 72-hour vitals chart became empty once the seeded data were more than
      three days old; the window now ends at each patient's latest reading.
\end{itemize}
